\chapter{Introduction}
\label{ch:introduction}

\section{Background and motivation}
\label{sec:motivation}

Job information in Germany and in the European Union is not in one place. It sits on public employment websites, commercial job boards, and company career pages. This is a problem for junior candidates, such as interns, working students, and early-career staff. It is also hard for people who need English-friendly jobs or visa information. Search takes a long time. Job texts mix German and English. Seniority words are not used in the same way. Visa rules are often missing or buried in long descriptions.

The OECD notes that artificial intelligence is changing labour markets. It affects how jobs are advertised, screened, and matched \cf{oecd2023ai}. Public services such as the Bundesagentur f\"ur Arbeit and EURES publish APIs and portals. However, a single candidate still needs to search many channels. Data products that unify public vacancy data can reduce search cost. They must still respect data quality, history, and European law.

DataForge is an applied data-science system that deals with this engineering problem. It already runs as a serverless lakehouse on Amazon Web Services in the Frankfurt region (\texttt{eu-central-1}). It collects data from several sources. It stores raw files in Bronze, historical job versions in Silver, and reports in Gold. It offers public APIs and a GitHub Pages website. This live system is the base of the thesis. The thesis does not start from an empty project.

Industry demand has moved from simple data collection to \emph{augmented analytics}. This means using machine learning and language processing to add labels, find similar jobs, and explain results \cf{provost2013ds,vial2019digital,demiroglu2015dsr}. In software companies, similar ideas appear under names such as hyperautomation and experiment-driven development \cf{fagerholm2017right,gartner2020hyper}. The research problem is not ``build a chatbot''. The problem is: how can generative AI be added to a live pipeline so that data history, repeatable tests, cost, and the law stay under control? Claims must also be possible to test.

\section{Problem statement}
\label{sec:problem}

When companies add large language models to existing data platforms, three problems appear often.

First, \emph{broken data history}. The model may overwrite job IDs or old versions. Then nobody can check what changed. Kimball and Ross warn that slowly changing dimensions need explicit rules. Type~2 history must not be replaced by silent rewrites \cf[p.~~220]{kimball2013dwt}. Second, \emph{quality that is not measured}. Demos show almost perfect ranking scores on tiny examples. They often skip a simple non-AI baseline. IR research has known for decades that ranking must be tested against labels \cf{manning2008ir,jarvelin2002ndcg}. Third, \emph{cost with no limit}. Expensive models are run on all documents. There is no sample rate, no budget, and no stop switch. These problems are part of a wider software issue: hidden extra work in machine-learning systems \cf{sculley2015debt,kreuzberger2023mlops}.

This study uses DataForge as one detailed case. The money limit is small (about \EUR{0}--25 per month in normal use, and a hard thesis cap of \EUR{40}). All cloud work stays in \texttt{eu-central-1}. EU data location is treated as a main design choice, not as a later extra.

\section{Stakeholders and intended use}
\label{sec:stakeholders}

Three stakeholder groups matter for this thesis.

\emph{Job seekers} want shorter lists of relevant vacancies. They care about language, seniority fit, and visa hints. \emph{Platform engineers} want a pipeline that they can operate, pause, and audit. \emph{Supervisors and examiners} want a thesis that connects theory, a built system, and honest evaluation. DataForge Match is designed as \emph{decision support} for seekers. It is not an employer screening tool. This boundary affects both metrics and law, as Chapters~\ref{ch:theory-ai} and~\ref{ch:discussion} explain.

\section{Aim and contribution of the study}
\label{sec:aim}

The aim is to design, build, and test a low-cost generative-AI layer on the existing DataForge lakehouse. The layer has a multi-provider model gateway, rule-based enrichment first, hybrid search for resume--job matching, a Match API that removes personal data, and a repeatable test set.

The thesis has four contributions. First, a reference architecture. Generative models add extra columns. They do not replace job history. Second, a matching experiment. It compares meaning-based search, keyword search, hybrid search, and the old rule-based wizard on a labelled set. Third, a simple cost model. It links model prices and sample rates to product cost. Fourth, a responsible-AI discussion. It maps an advice tool for job seekers to the GDPR and the EU Artificial Intelligence Act \cf{eu2016gdpr,eu2024aiact,veale2021demystifying}.

\section{Link to the Master Colloquium topics}
\label{sec:alignment}

This is an applied Data Science thesis. It is not only a software-engineering paper. It still matches several topics from the Master Colloquium of Prof.\ Dr.\ Rand Kouatly (Summer Semester 2026). The main topics are \emph{augmented analytics} and \emph{hyperautomation}. AI is used to add labels to job text and to automate matching that used to be keyword-based. Secondary topics are \emph{digitalisation of software-intensive services} (what changes when a live data product gets a generative layer) and \emph{automation of data cleaning} (rules first, then expensive language-model calls). Continuous experimentation is the test style: product claims are checked with held-out metrics, not only with a demo story \cf{fagerholm2017right}. Quantum computing, module clustering, and automated fault finding are nearby software topics. They are outside the scope of this thesis.

\section{Structure of the thesis}
\label{sec:structure}

Chapter~\ref{ch:theory-data} reviews digitalisation, lakehouse design, historical data modelling, data quality, and classical information retrieval. Chapter~\ref{ch:theory-ai} reviews language models, dense and hybrid search, retrieval-augmented generation, LLM operations, agents, and EU law. It then states the research gap. Chapter~\ref{ch:rqs} states the research questions and hypotheses. Chapter~\ref{ch:method} explains why design science and a single case study are used. Chapter~\ref{ch:results} reports what was built and what was measured. Chapter~\ref{ch:discussion} discusses the meaning of the results and the limits. Chapter~\ref{ch:implications} discusses theory, practice, and future research.

\section{Scope note for the reader}
\label{sec:scope-note}

The empirical core of this thesis is European public job text and a student-built AWS lakehouse. Results should be read as evidence for frugal GenAI integration under those constraints. They should not be read as a universal ranking of all foundation models, or as a labour-market policy evaluation. Where AWS quota blocked a live Bedrock cell, the text says so. Where labels are author-made, the text says so. That honesty is part of the scientific claim.
